% !PRINTED 656
\weberpage{656}

If we assume $n = p$, then $\mathfrak{p}_{p\xi} = \mathfrak{p}_\xi$ (\S.\,172), and hence also $k_{p\xi} = k_\xi$, and from (7) and (8) it follows that:
\weq{9}{k(p-1) \equiv 0 \ (\mathrm{mod}\ m).}

From this there now follow important consequences: If, in the first place, $p - 1$ is not divisible by $q$, it follows that $k$ must be divisible by $m$, and we therefore have the theorem:

4. \gesperrt{The prime factors of a prime $p$ which is not congruent to 1 with respect to the modulus $q$ are contained in $\omega$ only in such powers whose exponents are divisible by $m$.}

If, in the case where $m$ is a power of 2, $p - 1$ is divisible by 2 but not by 4, that is, if $p \equiv -1 \ (\mathrm{mod}\ 4)$, then it follows from (9) only that $k$ is divisible by $\tfrac{1}{2} m$, but not that $k$ is divisible by $m$. The exponent of the power of $\mathfrak{p}_\xi$ contained in $\omega$ is, apart from multiples of $m$, equal to $k$, and if $k$ is divisible by $m$, then everything is as in case 4. But if $k$ is divisible only by $\tfrac{1}{2} m$ and not by $m$, so that
\[
k \equiv \tfrac{1}{2} m \ (\mathrm{mod}\ m),
\]
then $k - \tfrac{1}{2} m$ is divisible by $m$, and, having regard to the decomposition (3), we can formulate the theorem thus:

5. \gesperrt{If $m$ is a power of 2 and $p$ a prime congruent to $-1$ with respect to the modulus 4, then the exponent $h$ can be so determined that all the prime factors of $p$ are contained in}
\[
\omega \sqrt{p}^{-hm}
\]
\gesperrt{only to such powers whose exponent is divisible by $m$.}

Finally, let $p - 1$ be divisible by $q$, or, if $m$ is even, by 4. If $m_1$ is the greatest common divisor of $m$ and $p - 1$, and if
\[
m = m_1 m_2,
\]
then by (9) $k$ must be divisible by $m_2$, and if we set
\[
k = h m_2
\]
then by (8) the following product is contained in $\omega$:
\weq{10}{\left(\prod^\xi \mathfrak{p}_\xi^{\xi'}\right)^{h m_2}.}
Moreover, the prime factors $\mathfrak{p}_\xi$ occur in $\omega$ only in such powers whose exponents are divisible by $m$.
